\begin{thebibliography}{99}

\bibitem{F18}
S. Fischler,
\emph{Shidlovsky's multiplicity estimate and irrationality of zeta values},
J. Aust. Math. Soc. \textbf{105} (2018), no.~2, 145--172.
\href{https://doi.org/10.1017/S1446788717000386}
{doi:10.1017/S1446788717000386}.

\bibitem{F20}
S. Fischler,
\emph{Irrationality of values of $L$-functions of Dirichlet characters},
J. Lond. Math. Soc. \textbf{101} (2020), no.~2, 857--876.
\href{https://doi.org/10.1112/jlms.12290}{doi:10.1112/jlms.12290}.
We use the \href{https://www.imo.universite-paris-saclay.fr/~stephane.fischler/Lschi.pdf}
{author revision of 10 October 2019}.

\bibitem{H12}
M. Huttner,
\emph{Riemann $P$-scheme, monodromy and diophantine approximations},
Indag. Math. \textbf{23} (2012), 522--546.
\href{https://doi.org/10.1016/j.indag.2012.03.002}
{doi:10.1016/j.indag.2012.03.002}.

\bibitem{K26}
S. Krisman,
\emph{One of $\beta(4),\beta(6),\ldots,\beta(36)$ is irrational},
preprint, 2026.
\href{https://hexagonmath.org/abs/2610.00285v2}{Hexagon 2610.00285v2}.

\bibitem{KT26}
V. Kumar and R. Tosi,
\emph{Irrationality of values of $L$-functions for real Dirichlet characters},
preprint, 30 October 2026.
\href{https://arxiv.org/abs/2610.41458v1}{arXiv:2610.41458v1}.

\bibitem{LS25}
L. Lai and J. Sprang,
\emph{Many $p$-adic odd zeta values are irrational},
preprint, version of 17 February 2025.
\href{https://arxiv.org/abs/2306.10393v2}{arXiv:2306.10393v2}.

\bibitem{LY20}
L. Lai and P. Yu,
\emph{A note on the number of irrational odd zeta values},
Compos. Math. \textbf{156} (2020), no.~8, 1699--1717.
\href{https://arxiv.org/abs/1911.08458v2}{arXiv:1911.08458v2}.

\bibitem{LZ22}
L. Lai and L. Zhou,
\emph{At least two of $\zeta(5),\zeta(7),\ldots,\zeta(35)$ are irrational},
\href{https://arxiv.org/abs/2103.00904v2}{arXiv:2103.00904v2},
Section~6, Theorem~6.1.

\bibitem{N11}
M. Nishimoto,
\emph{On the linear independence of the special values of a Dirichlet
series with periodic coefficients}, 2011.
\href{https://arxiv.org/abs/1102.3247v1}{arXiv:1102.3247v1}.

\bibitem{OAI26}
OpenAI,
\emph{Catalan's constant is irrational},
preprint, 24 September 2026.
\href{https://github.com/openai/math/blob/4eaee87548250b2682affda7d05d947c05554bae/preprints/Catalans-constant-is-irrational-September-24-2026/paper.pdf}
{OpenAI Math}.

\bibitem{RZ03}
T. Rivoal and W. Zudilin,
\emph{Diophantine properties of numbers related to Catalan's constant},
Math. Ann. \textbf{326} (2003), 705--721.
\href{https://doi.org/10.1007/s00208-003-0420-2}
{doi:10.1007/s00208-003-0420-2}.

\bibitem{S07}
V. N. Sorokin,
\emph{On the Zudilin--Rivoal theorem},
Mat. Zametki \textbf{81} (2007), no.~6, 912--923;
English translation: Math. Notes \textbf{81} (2007), no.~6, 817--826.
\href{https://doi.org/10.4213/mzm3742}{doi:10.4213/mzm3742}.

\bibitem{T26}
twinheterodox,
\emph{Linear independence of even Dirichlet beta values},
2026.
\href{https://github.com/twinheterodox/beta4/blob/095d309ae168757dbfa4063484d28186f7f41beb/even-beta-independence/output/pdf/even-beta-independence.pdf}
{Author version of 3 November 2026}.

\bibitem{Z19}
W. Zudilin,
\emph{Arithmetic of Catalan's constant and its relatives},
Abh. Math. Semin. Univ. Hambg. \textbf{89} (2019), 45--53.
\href{https://arxiv.org/abs/1804.09922v2}{arXiv:1804.09922v2}.

\end{thebibliography}
